\section{Trace class for the actual convolution-after-projection operator}
\label{sec:actual-lift-trace}

This section repairs a separate operator-order issue in the real
weighted model. It uses the actual quotient theorem MT-M26, with
its uniform interior containing-strip hypothesis. The new ingredient
is the differential form of translation on the specified weighted
$L^2$ space, including the finite boundary contribution from the
cusp of the weight at zero.

Fix $\alpha>0$, $s\ge0$, and write
\[
 w(x)=e^{2\alpha|x|}(1+x^2)^{s/2},\qquad
 H=L^2(\mathbb R,w(x)\,dx),\qquad (U_tf)(x)=f(x-t).
 \tag{MA-LIFT.1}
\]
Let $J$ be the Hilbert closure of the centered theta ideal, let
$Q=H/J$, and identify $Q$ isometrically with $Z=J^\perp$ by its
orthogonal representative. Write $P_Z$ and $P_G=1-P_Z$ for the
orthogonal projections. For $g\in C_c^\infty(\mathbb R)$, put
$C_gf=\int g(t)U_tf\,dt$. A multiplicative test $h$ corresponds
to $g(t)=h(e^t)$. All results below are for one fixed character
sector; no summation over varying conductors is inferred.

\paragraph{Input and zero-location scope.}
Assume precisely the hypothesis of MT-M26: a fixed closed strip
$[c_0,c_1]$ contains the real parts of all zeros of the relevant
completed function, and
\[
 [c_0-1/2,c_1-1/2]\subset(-\alpha,\alpha).
 \tag{MA-LIFT.2}
\]
For rational zeta one can take $[c_0,c_1]=[0,1]$ unconditionally,
so $\alpha>1/2$ suffices. Under the adopted Dirichlet input $A$,
one can take $c_0=c_1=1/2$ in those sectors, so every $\alpha>0$
suffices. For general Hecke characters this latter instantiation
uses the conditional input $M$. No unconditional partial-strip
version of MT-M26 is assumed.

\paragraph{MA-LIFT-001 (a trace-class resolvent power on the quotient).}
Under (MA-LIFT.2), let $A_Q$ generate the induced translation group
on $Q$. For a sufficiently large real $\kappa$ and an integer $k\ge2$,
\[
                  (\kappa-A_Q)^{-k}\in\mathcal S_1(Q).
 \tag{MA-LIFT.3}
\]
One may choose $k>B+2$, where $B$ is the polynomial exponent in
MT-M26's fixed-strip lifting bound (M26.9).

\emph{Proof.}
MT-M26 supplies finite block algebras $J_n$, uniformly bounded
maps $r_n:Q\to J_n$, polynomially bounded maps $s_n:J_n\to Q$, and
\[
 r_m s_n=\delta_{mn}\operatorname{id},\qquad
 d_n=\dim J_n=O(\log(2+|n|)),\qquad
 \|s_n\|\le C(2+|n|)^B.
 \tag{MA-LIFT.4}
\]
Their images span a dense subspace of $Q$: the smooth quotient has
absolute block synthesis, maps densely into $Q$, and its block
series converges in $Q$ by the polynomial lift bound. On
$s_n(J_n)$, translation is multiplication by the block class of
$e^{(z-1/2)t}$, by the smooth action and the exact zero-jet kernel
statement (M26.3).

Translation satisfies
\[
 \|U_t\|\le 2^{s/4}e^{\alpha|t|}(1+t^2)^{s/4}.
 \tag{MA-LIFT.5}
\]
The same holds on $Q$. Choose $\kappa>\alpha+1$; its excess
exponential decay absorbs the polynomial factor. The resolvent
power is the operator-norm convergent Laplace integral
\[
 (\kappa-A_Q)^{-k}=\frac1{(k-1)!}
 \int_0^\infty e^{-\kappa t}t^{k-1}U_t^Q\,dt.
 \tag{MA-LIFT.6}
\]
On block $n$ this is multiplication by
$F_k(z)=(\kappa-(z-1/2))^{-k}$. The contours used in MT-M26
lie in a fixed vertical strip, at heights comparable to $n$.
The quotient algebra norm thus gives
$\|M_{F_k,n}\|\le C_k(2+|n|)^{-k}$, and
\[
 \sum_n\|s_nM_{F_k,n}r_n\|_1
 \le C_k\sum_n d_n(2+|n|)^{B-k}<\infty\qquad(k>B+2).
 \tag{MA-LIFT.7}
\]
This trace-class sum agrees with (MA-LIFT.6) on each finite
block, hence on their dense span and then on $Q$.
\hfill$\square$

\paragraph{MA-LIFT-002 (the specified orthogonal lift is trace class).}
Under (MA-LIFT.2), the actual operator $S_g=C_gP_Z:H\to H$
is trace class for every $g\in C_c^\infty(\mathbb R)$, and
\[
 \operatorname{Tr}_H(C_gP_Z)
 =\operatorname{Tr}_Q(C_g^Q)
 =\sum_\rho m_\rho\,\mathcal Mg(\rho-1/2).
 \tag{MA-LIFT.8}
\]
The sum is absolutely convergent. On every fixed compact
test-support space, $g\mapsto S_g$ is continuous into trace norm.

\emph{Proof.}
The adjoint translation group is explicitly
\[
 (U_t^*u)(x)=\frac{w(x+t)}{w(x)}u(x+t),\qquad
 Lu=w^{-1}\partial_x(wu).
 \tag{MA-LIFT.9}
\]
It is strongly continuous by (MA-LIFT.5) and compactly supported
smooth approximation; $L$ is its generator. Derivatives in this
formula are weak derivatives. Multiplication by $w$ is a unitary
map from $H$ onto $L^2(w^{-1}dx)$, intertwining $U_t^*$ with
$v(x)\mapsto v(x+t)$. Consequently $u\in\mathcal D(L^k)$ exactly
when $v=wu$ has weak derivatives through order $k$ in
$L^2(w^{-1}dx)$, and $L^ju=w^{-1}\partial^jv$.
For the first derivative this follows from
$v(\cdot+t)-v=\int_0^t v'(\cdot+r)\,dr$ and strong continuity
of translation; the converse follows by testing a generator
difference quotient against compact smooth functions. Iteration
gives the asserted power domains.

Since $J$ is invariant under every $U_t$, $Z$ is invariant under
every $U_t^*$. This restricted group is the Hilbert adjoint of
the quotient group. MA-LIFT-001 gives
\[
 R_Z^k:=(\kappa-L|_Z)^{-k}\in\mathcal S_1(Z).
 \tag{MA-LIFT.10}
\]
Here $L|_Z$ means the generator of the restricted group. Its
resolvent is the restriction of $R=(\kappa-L)^{-1}$, by the
Laplace integral. The extended map
$\mathcal R_k=R^kP_Z:H\to H$ is therefore trace class and takes
values in $\mathcal D(L^k)$. It is bounded in the graph norm:
$L^jR^k$ is bounded for $j\le k$, using $LR=\kappa R-1$.

We now use the specified weight. On each open half-line it is
smooth, and
\[
 a_j(y):=w(y)\partial_y^j(w(y)^{-1})\quad(y\ne0)
       \ \in L^\infty(\mathbb R)\qquad(j\ge0).
 \tag{MA-LIFT.11}
\]
Indeed, $\log w=\pm2\alpha y+(s/2)\log(1+y^2)$ there; all
positive-order derivatives of this expression are bounded,
and each $a_j$ is a polynomial in them. All one-sided
derivatives of $w^{-1}$ at zero also exist.

For $u\in\mathcal D(L^k)$ put $v=wu$. Then
\[
 C_g(\kappa-L)^ku(x)=
 \int_{\mathbb R}\frac{g(x-y)}{w(y)}
                    (\kappa-\partial_y)^kv(y)\,dy.
 \tag{MA-LIFT.12}
\]
Integrate by parts separately on $y<0$ and $y>0$. Terms at
infinity vanish because $g(x-y)$ has compact support in $y$.
The bulk term is
\[
 B_{g,k}u(x)=\int_{y\ne0}
 w(y)(\kappa+\partial_y)^k
           \left(\frac{g(x-y)}{w(y)}\right)u(y)\,dy.
 \tag{MA-LIFT.13}
\]
Expansion makes this a finite sum of $C_{g^{(i)}}M_{a_j}$
with coefficients depending on $\kappa,k$. Each summand is
bounded on $H$ by (MA-LIFT.5), the weighted $L^1$ norm of
$g^{(i)}$, and (MA-LIFT.11). Hence $B_{g,k}$ is bounded.

The remaining terms have the form
\[
 F_{g,k}u=\sum_{m=0}^{k-1}\phi_{g,k,m}\,\ell_m(u),
 \qquad \ell_m(u)=\partial^m(wu)(0).
 \tag{MA-LIFT.14}
\]
Each $\phi_{g,k,m}$ is a finite linear combination of
$g,g',\ldots,g^{(k-1)}$ with coefficients given by jumps of
one-sided derivatives of $w^{-1}$ at zero. In particular
$\phi_{g,k,m}\in C_c^\infty(\mathbb R)\subset H$.
Explicitly the boundary terms of the derivative of order $j$
are linear combinations of
\[
 \bigl[\partial_y^l(g(x-y)/w(y))\bigr]_{0-}^{0+}
                     v^{(j-1-l)}(0),\qquad 0\le l<j\le k.
\]
These common derivatives of $v$ exist since $v\in H^k_{\rm loc}$.
This is the ordinary Sobolev integration-by-parts formula;
the cusp $|x|$ has not been treated as globally smooth.

The one-dimensional Sobolev trace inequality gives
\[
 |\ell_m(u)|\le C_k\|wu\|_{H^k([-1,1])}
 \le C_{k,\alpha,s}\sum_{j=0}^k\|L^ju\|_H
 \qquad(0\le m\le k-1).
 \tag{MA-LIFT.15}
\]
For the last inequality use boundedness above and below of $w$
on this interval. Thus $\ell_mR^kP_Z$ is bounded on $H$.
The trace inequality needed here is elementary: the fundamental
theorem of calculus, followed by averaging and Cauchy--Schwarz,
bounds $|q(0)|$ by $C(\|q\|_2+\|q'\|_2)$ on $[-1,1]$.
Apply it to $q=v^{(m)}$, with $m+1\le k$.

Apply (MA-LIFT.12)--(MA-LIFT.14) to $u=\mathcal R_kf$.
Since $(\kappa-L)^ku=P_Zf$, the exact identity is
\[
 C_gP_Z=B_{g,k}\mathcal R_k+
     \sum_{m=0}^{k-1}\phi_{g,k,m}\otimes(\ell_m\mathcal R_k).
 \tag{MA-LIFT.16}
\]
The first term is trace class and the rest have finite rank.
More quantitatively,
\[
 \|C_gP_Z\|_1\le \|B_{g,k}\|\,\|R_Z^k\|_1+
   \sum_{m=0}^{k-1}\|\phi_{g,k,m}\|_H
                    \|\ell_mR^kP_Z\|_{H'}.
 \tag{MA-LIFT.17}
\]
All dependence on $g$ is controlled by finitely many derivatives
on its fixed compact support, proving the stated continuity.

Finally, trace cyclicity with the bounded projection gives
$\operatorname{Tr}(C_gP_Z)=\operatorname{Tr}(P_ZC_gP_Z)$.
The latter is the orthogonal representative of quotient
convolution, with zero action on $J$. MT-M26 computes its
trace as (MA-LIFT.8), with absolute convergence.
\hfill$\square$

\paragraph{MA-LIFT-003 (actual cutoff spectral and geometric traces).}
Under the same hypotheses, let $\rho_\nu\in C_c^\infty(\mathbb R)$
have uniformly bounded suprema and tend to $1$ almost everywhere.
Then
\[
       \|(M_{\rho_\nu}-1)C_gP_Z\|_1\longrightarrow0.
 \tag{MA-LIFT.18}
\]
For every compactly supported smooth $\rho$, $M_\rho C_g$
is trace class on $H$ and
\[
 \operatorname{Tr}_H(M_\rho C_g)=g(0)\int_{\mathbb R}\rho(x)\,dx.
 \tag{MA-LIFT.19}
\]
In particular, if $g(0)=0$, the actual geometric-order operator
$G_{\rho,g}=M_\rho C_gP_G$ is trace class, and
\[
 \lim_\nu\operatorname{Tr}_H G_{\rho_\nu,g}
       =-\sum_\rho m_\rho\,\mathcal Mg(\rho-1/2).
 \tag{MA-LIFT.20}
\]
Here $\rho$ in the sum denotes zeros, independently of the
cutoff notation.

\emph{Proof.}
Dominated convergence gives $M_{\rho_\nu}\to1$ strongly on $H$,
with uniformly bounded operator norms. Strong convergence of a
uniformly bounded family, followed by a fixed trace-class
operator, is convergence in trace norm: first verify it for
finite-rank operators and then approximate in trace norm.
MA-LIFT-002 therefore gives (MA-LIFT.18).

On ordinary $L^2(\mathbb R)$ the kernel $\rho(x)g(x-y)$ is
smooth and compactly supported in both variables, so it is
trace class with trace $g(0)\int\rho$. An absolutely summable
Hermite expansion proves this as in MA-LA-003.
Conjugate to $H$ by multiplication by $\sqrt w$. Choose a smooth
compact $\eta$ equal to one on both coordinate projections of
the kernel support. The conjugated operator is
\[
 M_{\eta\sqrt w}\,(M_\rho C_g)_{L^2}\,M_{\eta/\sqrt w}.
\]
Both outside factors are bounded. Trace cyclicity replaces their
product by $M_{\eta^2}$, which leaves the middle kernel unchanged.
This proves (MA-LIFT.19), including trace class, without requiring
smoothness of the weight at zero.

Now $G_{\rho,g}=M_\rho C_g-M_\rho C_gP_Z$ is a difference of
trace-class operators. If $g(0)=0$, its trace equals
$-\operatorname{Tr}(M_\rho C_gP_Z)$ by (MA-LIFT.19).
Use (MA-LIFT.18) and (MA-LIFT.8) to pass to the limit.
\hfill$\square$

\paragraph{What this closes, and what it does not.}
MA-LIFT-002 replaces the unfinished lifted trace-class estimate
in Masullo, p.~464, equations~(10.4.2.25)--(10.4.2.34), for this
specified fixed-sector weighted model under (MA-LIFT.2).
MA-LIFT-003 supplies cutoff convergence in the operator order
of Theorem~10.4.2.8. It discharges the actual lifted trace-class
premise of F-SM-18/F-SM-21 in the rational wide-strip case, and
in every rational narrow strip when $A$ is explicitly used.
It does not contradict the abstract obstruction MA-TR-102.

The arithmetic quotient is used only through the verified
finite block data and the uniform interior strip hypothesis.
The rest is a calculation with the specified weight and actual
translations. Equation (MA-LIFT.20) is a trace theorem on that
actual global Hilbert space, rather than a scalar assigned from
the explicit formula. Nevertheless its computation is spectral.
Identifying it with an independently transported local geometric
kernel carrying the prime terms remains the global-to-local
question. Using $A$ or $M$ to supply (MA-LIFT.2) in every narrow
strip and then presenting this as an independent proof of their
zero-location conclusion would be circular.

